\chapter*{Conclusiones}
\addcontentsline{toc}{chapter}{Conclusiones}

% Las conclusiones tampoco cuentan como capítulo, ctm

\noindent En esta tesis, me propongo evaluar el impacto causal de las cámaras de velocidad en el número de incidentes viales en la Ciudad de México. Para lograr esto, evalúo el impacto de la transición del programa de \textit{Fotomultas} al sistema de \textit{Fotocívicas} en la Ciudad de México a partir de abril de 2019. El cambio no fue menor: implicó una redistribución de los radares de velocidad hacia zonas de alta siniestralidad y, sobre todo, una transformación en la lógica de sanción, priorizando el resarcimiento cívico y educativo sobre el fin meramente recaudatorio. Mi hipótesis inicial, como lo sugería la literatura, era que este despliegue generaría una reducción localizada de los incidentes viales, con especial énfasis en los eventos de mayor severidad a una distancia de 300 metros alrededor de los dispositivos. 

Para poder analizar el impacto, hago uso de dos estrategias complementarias: el \textit{Propensity Score Matching (PSM)}, para construir un grupo de control comparable, y los modelos de \textit{Diferencia en Diferencias}, para realizar la estimación de los efectos. Sin embargo, los resultados obtenidos tras este proceso contradicen la expectativa de un impacto puntual. Al comparar el grupo de control construido para ser estadísticamente similar --controlando por historial de accidentes, infraestructura vial y proxies de movilidad--, no se encontró evidencia estadística de que las cámaras generen una reducción localizada, ni en el número total de incidentes, ni en la tasa de incidentes en sus áreas de influencia. 

Más allá de estos hallazgos, la tesis deja cuatro aportes para la evaluación de políticas de control de velocidad en la Ciudad de México. El primero es una unidad espacial de análisis alineada con el carácter local del tratamiento, al estudiar radios centrados en cada cámara en lugar de áreas agregadas. El segundo es un contrafactual más creíble, construido mediante emparejamiento sobre siniestralidad previa, infraestructura vial y medidas de exposición. El tercero es una estrategia de medición que integra registros de la afluencia del Metro y estaciones permanentes de conteo de la SICT para incorporar cambios exógenos en movilidad durante la pandemia. El cuarto es interpretativo: distinguir entre efectos agregados del programa y efectos localizados de los radares permite visitar con mayor precisión la evidencia previa para la ciudad. 

Es importante destacar que esta ausencia de significancia en el efecto localizado no implica que la siniestralidad no haya disminuido. Por el contrario, algo que sí se pudo identificar fue que la variable que captura los cambios en el tiempo, mostró una reducción generalizada de la siniestralidad tanto en las áreas con una cámara de velocidad como en aquellas que no tienen una, siendo esta reducción de entre el 15\% y el 20\% tras la implementación del programa. Este dato, que coincide con las evaluaciones previas de la SEMOVI (2020), Quintero et. al (2022), sugiere que estos efectos no están necesariamente relacionados con la presencia física de un radar de velocidad, sino que más bien fueron reducciones generalizadas. 

Derivado de esto, surgen dos hipótesis principales que podrían explicar la coexistencia de una mejora sistémica frente a la ausencia de un efecto localizado. La primera reside en la posibilidad de que el programa de \textit{Fotocívicas} haya logrado efectivamente un cambio de comportamiento generalizado en la cultura vial de la Ciudad de México. Bajo esta lógica, la percepción de vigilancia y el carácter público de la ubicación de los radares habrían influido en la conducción de manera global, mitigando riesgos más allá de los puntos específicos de monitoreo. La segunda hipótesis sugiere un fenómeno de desplazamiento espacial: es posible que los conductores hayan modificado sus trayectorias habituales para evitar no solo las zonas con radares, sino también aquellas áreas con configuraciones viales similares donde se percibe una mayor probabilidad de sanción.

Lamentablemente, la falta de datos granulares sobre la afluencia vehicular a nivel de intersección impide validar de manera definitiva estas conjeturas. Sin una métrica precisa de flujo por tramo, no es posible distinguir si la reducción observada en los incidentes responde a una conducción más segura o simplemente a una redistribución del tráfico hacia vías no monitoreadas. Por lo tanto, con la evidencia disponible, mi análisis solo puede concluir que las cámaras de velocidad no presentan un efecto localizado estadísticamente significativo sobre el número de incidentes viales ni sobre su severidad en los radios estudiados.

A esta limitación se suma la complejidad de un esquema de sanciones que no es homogéneo. El hecho de que una proporción considerable de las infracciones siga teniendo un carácter económico --para vehículos con placas de otras entidades federativas-- introduce una variable de sesgo adicional. Esta dualidad en los incentivos podría estar diluyendo el impacto esperado de un sistema diseñado originalmente para el resarcimiento cívico, dificultando la identificación de si el resultado es atribuible a los radares o a la respuesta diferenciada de los usuarios ante el tipo de multa.

En última instancia, contar con datos que permitan medir con precisión el flujo vehicular en cada intersección será fundamental para futuras investigaciones. Esta información no solo permitiría robustecer la identificación del efecto causal, sino que también sentaría las bases para diseñar intervenciones más eficaces y propuestas de política pública que logren integrar de manera estratégica nuevos radares en el sistema de seguridad vial de la ciudad.